\documentclass[11pt,a4paper]{article}

\usepackage[margin=2.4cm]{geometry}
\usepackage{amsmath,amssymb}
\usepackage{booktabs}
\usepackage{amsthm}
\usepackage{siunitx}
\usepackage[hidelinks]{hyperref}

\sisetup{group-separator={,},group-minimum-digits=4}
\DeclareSIUnit{\rupee}{Rs}
\DeclareSIUnit{\USD}{\$}
\DeclareSIUnit{\yr}{yr}

\newtheorem{proposition}{Proposition}

\newcommand{\Bmax}{B^{\max}}
\newcommand{\dQ}{\Delta Q}
\newcommand{\etain}{\eta^{\mathrm{in}}}
\newcommand{\etaout}{\eta^{\mathrm{out}}}

\title{Battery Health in the Customer Menu:\\
Session Capacity Fade and its Monetary Equivalent}
\author{}
\date{}

\begin{document}
\maketitle

\begin{abstract}
\noindent
This note states the battery-health component of the residential EV
charging menu: the semi-empirical capacity-fade model evaluated for one
connected session, its bidirectional (V2G) extension, and the mapping
from fade to rupees that makes health commensurable with the electricity
bill. Equations are numbered independently; where they correspond to the
NPSC manuscript this is noted. Every numerical claim in
Section~\ref{sec:verification} was re-derived from the parameter table
alone and agrees with the implementation to \num{1.4e-20}. Parameter
provenance and its limits are stated in Section~\ref{sec:provenance};
these are literature-anchor calibrations, not fitted coefficients, and
the monetary conversion rests on a single exogenous assumption.
\end{abstract}

\section{Session fade}

Let a connected session span the interval set $\mathcal{T} =
\{0,\dots,n-1\}$ with durations $\Delta t_k$ (hours), battery energies
$E_k$ at the $n+1$ state boundaries, and state of charge $s_k = E_k /
\Bmax$. Capacity fade over the session decomposes into three terms:
%
\begin{equation}
\dQ \;=\; \dQ^{\mathrm{cal,win}} \;+\; \dQ^{\mathrm{cal,park}}
       \;+\; \dQ^{\mathrm{cyc}} .
\label{eq:total}
\end{equation}
%
All three are dimensionless fractions of nominal usable capacity. This
replaces the normalised $0$--$100$ health score $H$ of the manuscript;
\eqref{eq:total} is a \emph{damage} measure, so the dominance direction
reverses, $\dQ_a \le \dQ_b$ where the manuscript wrote $H_a \ge H_b$.

\subsection{Calendar stress}

Calendar ageing is driven by a function of the state of charge,
%
\begin{equation}
g(s) \;=\; a_0 \;+\; a_1 s \;+\; a_2 \bigl(\max\{s - s^{\mathrm{knee}},\,0\}\bigr)^{2}
          \;+\; a_3 \bigl(s - s_c\bigr)^{3},
\label{eq:gstress}
\end{equation}
%
in fractional capacity loss per year at the reference temperature and
anchor age. Two shapes are admitted. A hinge at $s^{\mathrm{knee}}$
encodes accelerated degradation above a threshold and is flat beneath
it. A cubic centred on $s_c$ rises monotonically with a plateau at the
centre. They are not interchangeable, and which one is used decides
whether discharging a pack at moderate state of charge carries any
calendar benefit.

\paragraph{LFP.} The coefficients are those of Naumann et al.\
(2018)~\cite{naumann2018}, converted to the differential form of
Section~\ref{sec:parameters}. That study reports a cubic,
$a_1 = a_2 = 0$ with $s_c = 0.5$, and calendar loss rising
monotonically with state of charge by a factor of \num{3.915} from
empty to full. An earlier version of this formulation used a hinge at
$s^{\mathrm{knee}} = 0.70$ with $a_1 = 0$, which is flat below
\SI{70}{\percent} and gives a full-to-empty spread of only
\num{1.24}; it is not what the cited study measured and has been
replaced.

Two multiplicative factors move $g$ off its reference point. Temperature
follows Arrhenius,
%
\begin{equation}
\theta(T) \;=\; \exp\!\left[-\frac{E_a}{R}
  \left(\frac{1}{T} - \frac{1}{T^{\mathrm{ref}}}\right)\right],
\qquad T,\,T^{\mathrm{ref}} \text{ in kelvin},
\label{eq:arrhenius}
\end{equation}
%
with $R = \SI{8.314462618}{\joule\per\mole\per\kelvin}$. Ageing is a
sublinear power law in calendar time, $Q \propto t^{\alpha}$ with
$\alpha \in (0,1]$, so the \emph{local} rate at pack age $t$ relative to
the anchor age $t_{0}$ is the ratio of derivatives,
%
\begin{equation}
\psi(t) \;=\;
  \frac{\alpha\, t^{\alpha-1}}{\alpha\, t_{0}^{\alpha-1}}
  \;=\; \left(\frac{t}{t_{0}}\right)^{\alpha-1}.
\label{eq:agefactor}
\end{equation}
%
Since $\alpha < 1$, $\psi$ decreases with age: an older pack fades more
slowly per unit time. The model is therefore incremental, charging each
session the marginal rate at the pack's current age rather than an
average, which is the correct treatment for a power law.

The anchor age $t_{0}$ is a normalisation, not a physical quantity. It
fixes the age at which $g(s)$ in \eqref{eq:gstress} is literally the
fractional loss per year, so that $\psi(t_{0}) = 1$ by construction.
Taking $t_{0} = \SI{1}{\yr}$ reduces \eqref{eq:agefactor} to
$t^{\alpha-1}$ with $t$ in years and places the coefficients near the
one-to-two-year span of the calendar-ageing campaigns they are drawn
from; the singularity at $t\to 0$ rules out anchoring at zero.
Consequently $t_{0}$ and $(a_0,a_1,a_2,a_3)$ are \emph{not} independent:
re-anchoring to $t_0'$ requires scaling every $a_i$ by
$(t_0'/t_{0})^{1-\alpha}$, failing which the model changes silently.
For $t_0' = \SI{2}{\yr}$ that factor is \num{1.189} (NMC) and
\num{1.414} (LFP).

\subsection{Charging-window and parked-day calendar fade}

Over the connected session, stress is accumulated at the left endpoint
of each interval:
%
\begin{equation}
\dQ^{\mathrm{cal,win}} \;=\; \theta(T)\,\psi(t)
  \sum_{k=0}^{n-1} g\!\left(\min\{\max\{s_k,0\},1\}\right)
  \frac{\Delta t_k}{H_{\mathrm{yr}}},
  \qquad H_{\mathrm{yr}} = \num{8760}.
\label{eq:calwin}
\end{equation}
%
The remainder of the day is spent parked at the delivered state of
charge, less the energy consumed by the trip $E^{\mathrm{trip}}$:
%
\begin{equation}
s^{\mathrm{park}} \;=\;
  \frac{\max\{\zeta,\,z\}\,\Bmax - E^{\mathrm{trip}}}{\Bmax},
\qquad
\dQ^{\mathrm{cal,park}} \;=\; g\!\left(s^{\mathrm{park}}\right)
  \theta(T)\,\psi(t)\,\frac{H_{\mathrm{park}}}{H_{\mathrm{yr}}},
\label{eq:calpark}
\end{equation}
%
where $z$ is the requested target, $\zeta$ the declared departure floor,
and $H_{\mathrm{park}}$ the parked hours. Because terminal energy is
pinned by equality to $\max\{\zeta,z\}\Bmax$,
\eqref{eq:calpark} is constant at fixed $(z,\zeta)$ and independent of
the charging schedule. This is Proposition~1: the parked term is a
constant of the request and may be excluded from the trajectory
objective without changing its argument minimum.

\subsection{Cycle fade}

Cycle ageing is a power law in accumulated equivalent full cycles $N$,
modulated by depth of discharge and rate. Its local slope is
%
\begin{equation}
\sigma(N) \;=\; \beta \,\max\{N,\,N^{\min}\}^{\,\beta-1},
\label{eq:cycslope}
\end{equation}
%
with $N^{\min}>0$ a regularisation floor for the singular derivative at
$N=0$. Then
%
\begin{equation}
\dQ^{\mathrm{cyc}} \;=\; c_{\mathrm{ref}}
  \bigl(1 + c_D\,\delta\bigr)
  \bigl(1 + c_C\,\hat{C}\bigr)\,
  \sigma(N)\;\varphi ,
\label{eq:cyc}
\end{equation}
%
where $\varphi$ is battery-side throughput as a fraction of capacity,
$\delta$ the state-of-charge excursion, and $\hat{C}$ the peak C-rate.
These three are where the bidirectional extension enters.

\section{Bidirectional extension}
\label{sec:v2g}

Write $E^{\mathrm{in}}_k \ge 0$ and $E^{\mathrm{out}}_k \ge 0$ for grid
import and export in interval $k$. Battery energy evolves as
%
\begin{equation}
E_{k+1} \;=\; E_k \;+\; \etain E^{\mathrm{in}}_k
  \;-\; \frac{E^{\mathrm{out}}_k}{\etaout}.
\label{eq:state}
\end{equation}
%
Throughput counts both directions on the battery side:
%
\begin{equation}
\varphi \;=\; \frac{1}{\Bmax}
  \left( \etain \sum_{k} E^{\mathrm{in}}_k
       + \frac{1}{\etaout}\sum_{k} E^{\mathrm{out}}_k \right),
\qquad
\Delta N \;=\; \frac{\varphi}{2}.
\label{eq:phi}
\end{equation}
%
The equivalent-full-cycle convention in \eqref{eq:phi} is a definition,
not an inheritance: one EFC is $2\Bmax$ of battery-side throughput, so a
full charge followed by a full discharge advances $N$ by exactly one.

Under charge-only operation the state of charge is monotone, and
throughput coincides with the state-of-charge swing, so $\varphi$ serves
as a depth proxy. Export breaks that, and depth must be stated
explicitly:
%
\begin{equation}
\delta \;=\; \max_{k} s_k \;-\; \min_{k} s_k .
\label{eq:delta}
\end{equation}
%
\begin{proposition}[backward compatibility]
If $E^{\mathrm{out}}_k \equiv 0$ and no trip occurs within the session,
then $\delta = \varphi$ exactly.
\end{proposition}
%
\noindent
Peak C-rate likewise takes the larger of the two directions:
%
\begin{equation}
\hat{C} \;=\; \frac{1}{\Bmax} \max_{k}
  \max\!\left\{
    \frac{\etain E^{\mathrm{in}}_k}{\Delta t_k},\;
    \lambda_d\,\frac{E^{\mathrm{out}}_k}{\etaout \Delta t_k}
  \right\}.
\label{eq:chat}
\end{equation}
%
Here $\lambda_d$ is a discharge-rate severity multiplier. It has no
anchor in the parameter sources used here, neither of which reports a
discharge-side C-rate coefficient, and is therefore exposed as a
sensitivity parameter rather than claimed as calibrated.

Two structural consequences follow. First, $\varphi$ becomes a decision
variable under export, so \eqref{eq:cyc} can no longer be treated as a
constant of the request and must re-enter the trajectory objective.
Second, $\bigl(1 + c_D\delta\bigr)\bigl(1 + c_C\hat{C}\bigr)\varphi$ is
not convex, so convexity claims valid for the charge-only problem do not
survive. A claim made in an earlier version of this formulation is
\emph{withdrawn} here. It held that, because LFP calendar ageing was
flat below the knee, export from beneath it carried no calendar
benefit. That rested on the hinge form; under the cubic actually
reported by Naumann et al.\ \cite{naumann2018}, $g$ is strictly
increasing in $s$ everywhere, so lowering the state of charge does
reduce calendar ageing at every operating point. The quantitative
claim is modest --- moving a pack from \SI{65}{\percent} to
\SI{45}{\percent} reduces $g$ by roughly a quarter --- and it is
offset by the cycle fade the export itself incurs, so no headline
claim that V2G improves battery health follows either way.

\section{Monetary equivalent}
\subsection{Pricing a session}

The menu offers four attributes to a customer who must trade battery
wear against the electricity bill. A normalised health index does not
support that trade; a price does. Since $\dQ$ is the \emph{fraction} of
usable capacity consumed, multiplying by capacity and pack cost gives
rupees of battery spent:
%
\begin{equation}
W \;=\; \dQ \,\cdot\, \Bmax \,\cdot\, \pi_{\mathrm{chem}},
\label{eq:wear}
\end{equation}
%
with $\pi_{\mathrm{chem}}$ the chemistry-specific pack cost in
\si{\rupee\per\kilo\watt\hour}. Wear and bill are then in the same unit,
and the only free quantity in the choice model is their ratio: how much
of a rupee of consumed pack a customer treats as a rupee out of pocket.

\subsection{What to display}

Equation~\eqref{eq:wear} prices the whole session, which is the correct
quantity for the choice model and the wrong one to put in front of a
customer. The reason is that a large part of $W$ is incurred whichever
offer is taken. The parked-day term \eqref{eq:calpark} depends only on
the delivered state of charge and the trip, so it accrues even when the
customer declines to charge altogether; in the worked example of
Section~\ref{sec:verification} it is the largest of the three
components, at \SI{43}{\percent} of total fade. A customer comparing two
rows of a menu is choosing between them, not between charging and
owning a battery, so the figure shown against an offer should be the
wear its selection \emph{adds}.

Let $\dQ_0$ be the smallest total fade among the offers of one menu. In
practice the minimiser is the decline-to-charge option, so $\dQ_0$ is
the fade the customer incurs by doing nothing. For an offer of fade
$\dQ$ and saving $S$, define
%
\begin{equation}
W^{\mathrm{marg}} \;=\; \bigl(\dQ - \dQ_0\bigr)\Bmax\,\pi_{\mathrm{chem}},
\qquad
\mathrm{Net} \;=\; S - W^{\mathrm{marg}}.
\label{eq:net}
\end{equation}
%
$W^{\mathrm{marg}}$ is the pack spend attributable to the decision and
$\mathrm{Net}$ is what the offer is worth once that spend is paid for.
Three properties of \eqref{eq:net} should travel with any figure
computed from it.

It is a presentation convention, not part of the model. The choice
model uses absolute $W$ from \eqref{eq:wear}. Within one request the two
differ by the constant $\dQ_0 \Bmax \pi_{\mathrm{chem}}$, which cancels
in a logit and so leaves choice probabilities unchanged; across requests
it does not cancel, and the two conventions must not be mixed.

$\dQ_0$ is a property of the menu, not of the vehicle. Regenerate the
menu with a different tier ladder and $\dQ_0$ moves, so
$\mathrm{Net}$ ranks offers within one menu and is not comparable
across menus or across sessions.

$\mathrm{Net}$ can be negative, and this is the point of computing it.
An offer whose saving is smaller than the pack it consumes destroys
value; in the worked menu of Section~\ref{sec:verification} two of
fifteen displayed offers are negative, which is not visible when the
health column reports a raw fade fraction.

\section{Parameters}
\label{sec:parameters}

\begin{table}[htbp]
\centering
\caption{Degradation parameters as implemented. LFP calendar
coefficients are Naumann et al.\ (2018)~\cite{naumann2018} Table~7,
converted to the per-year differential form by
$g(s) = \tfrac{1}{2}\,(k_{\mathrm{ref}}/100)\sqrt{S}\;k_{\mathrm{SOC}}(s)$
with $S = 8760 \times 3600$. Cycle coefficients for both chemistries,
and all NMC calendar coefficients, are not drawn from either cited
study; see Section~\ref{sec:provenance}. The upper block is
chemistry-specific and immutable; the lower block holds scenario
settings shared across chemistries, which is why no reference age or
reference temperature appears in a per-chemistry coefficient listing.}
\label{tab:params}
\begin{tabular}{llrr}
\toprule
Symbol & Quantity & LFP & NMC \\
\midrule
$a_0$              & calendar intercept            & \num{0.021258} & \num{0.0073} \\
$a_1$              & calendar linear term          & \num{0}      & \num{0.0210} \\
$a_2$              & calendar hinge coefficient    & \num{0}      & \num{0.0475} \\
$s^{\mathrm{knee}}$& calendar knee                 & ---          & \num{0.80}   \\
$a_3$              & calendar cubic coefficient    & \num{0.100862} & \num{0} \\
$s_c$              & cubic centre                  & \num{0.50}   & ---          \\
$\alpha$           & calendar time exponent        & \num{0.50}   & \num{0.75}   \\
$E_a$              & activation energy (\si{\joule\per\mole}) & \num{17126} & \num{24000} \\
$T^{\mathrm{ref}}$ & calendar reference temperature & \SI{25}{\celsius} & \SI{30}{\celsius} \\
$c_{\mathrm{ref}}$ & cycle reference coefficient   & \num{0.0007} & \num{0.0009} \\
$c_D$              & depth-of-discharge coefficient& \num{1.00}   & \num{1.00}   \\
$c_C$              & C-rate coefficient            & \num{1.00}   & \num{1.00}   \\
$\beta$            & cycle throughput exponent     & \num{0.50}   & \num{0.50}   \\
\midrule
$t_{0}$ & calendar anchor age            & \multicolumn{2}{c}{\SI{1}{\yr}} \\
$H_{\mathrm{park}}$& parked hours per day          & \multicolumn{2}{c}{\num{16}} \\
$N^{\min}$         & cycle-slope regularisation    & \multicolumn{2}{c}{\num{1}} \\
$\pi_{\mathrm{chem}}$ & pack cost (\si{\rupee\per\kilo\watt\hour}) & \num{7850} & \num{12400} \\
\bottomrule
\end{tabular}
\end{table}

\section{Worked example}
\label{sec:verification}

A \SI{78}{\kilo\watt\hour} NMC pack, three years old, \num{500}
accumulated EFC, at \SI{25}{\celsius}; connected \num{20}:\num{00} to
\num{07}:\num{00}, arriving at \SI{65}{\percent}. The factors are
%
\[
\theta = \num{0.852415218}, \qquad
\psi = 3^{-0.25} = \num{0.759835686}, \qquad
\sigma(500) = 0.5 \cdot 500^{-0.5} = \num{0.022360680}.
\]
%
For the highest-saving offer on the menu (target \SI{94}{\percent},
$S = \SI{143.32}{\rupee}$):

\begin{table}[htbp]
\centering
\begin{tabular}{lrr}
\toprule
Component & Value & Share \\
\midrule
$\dQ^{\mathrm{cal,win}}$  & \num{1.695736e-5} & \SI{24.4}{\percent} \\
$\dQ^{\mathrm{cal,park}}$ & \num{2.987358e-5} & \SI{43.0}{\percent} \\
$\dQ^{\mathrm{cyc}}$      & \num{2.270323e-5} & \SI{32.7}{\percent} \\
\midrule
$\dQ$ (total)             & \num{6.953417e-5} & \SI{0.00695}{\percent} of pack \\
\bottomrule
\end{tabular}
\end{table}

\noindent
Applying \eqref{eq:wear} and \eqref{eq:net} with $\dQ_0 =
\num{3.952973e-5}$:
%
\[
W = \num{6.953417e-5}\cdot 78 \cdot \num{12400} = \SI{67.25}{\rupee},
\qquad
W^{\mathrm{marg}} = \SI{29.02}{\rupee},
\qquad
\mathrm{Net} = \SI{114.30}{\rupee}.
\]

\paragraph{Verification.}
All \num{114} Pareto-efficient offers of this menu were re-derived from
Table~\ref{tab:params} alone, retyping \eqref{eq:gstress}--\eqref{eq:cyc}
rather than calling the implementation's own fade routines. Worst
absolute disagreement against the reported values: \num{1.016e-20} on
$\dQ^{\mathrm{cal,win}}$, exactly zero on $\dQ^{\mathrm{cal,park}}$ and
$\dQ^{\mathrm{cyc}}$, and \num{1.355e-20} on the total.

\paragraph{Sensitivity to the pack cost.}
$W^{\mathrm{marg}}$ and Net are linear in $\pi_{\mathrm{chem}}$:

\begin{table}[htbp]
\centering
\begin{tabular}{lrrr}
\toprule
Scenario & $\pi_{\mathrm{NMC}}$ (\si{\rupee\per\kilo\watt\hour}) & $W^{\mathrm{marg}}$ & Net \\
\midrule
$0.75\times$           & \num{9300}  & \SI{21.77}{\rupee} & \SI{121.55}{\rupee} \\
baseline               & \num{12400} & \SI{29.02}{\rupee} & \SI{114.30}{\rupee} \\
$1.25\times$           & \num{15500} & \SI{36.28}{\rupee} & \SI{107.04}{\rupee} \\
BEV-pack average       & \num{9584}  & \SI{22.43}{\rupee} & \SI{120.89}{\rupee} \\
\bottomrule
\end{tabular}
\end{table}

\section{Provenance and limitations}
\label{sec:provenance}

\paragraph{LFP calendar: reproduced, not approximated.}
The LFP calendar coefficients are Table~7 of Naumann et
al.~\cite{naumann2018}, converted as stated in
Table~\ref{tab:params}. The implementation reproduces that paper's own
published projection --- \SI{19.0}{\percent} capacity loss at 20 years,
\SI{25}{\celsius}, \SI{50}{\percent} SOC --- as \SI{19.01}{\percent},
and agrees with their Eq.~(3) to a relative \num{1e-6} across
temperatures \SIrange{0}{60}{\celsius}, all states of charge and
horizons from six months to twenty years. Both checks are asserted in
the test suite, so a later edit to a coefficient breaks a test rather
than quietly invalidating the citation.

\paragraph{NMC: Schmalstieg et al.\ (2014), in part.}
The NMC coefficients come from Schmalstieg et
al.~\cite{schmalstieg2014}, a Sanyo UR18650E study of more than sixty
cells covering both calendar and cycle ageing. Four quantities are
theirs. The calendar time exponent $\alpha = 0.75$ is their Section~3.2
choice, adopted after a $\sqrt{t}$ law fitted their data poorly. The
activation energy follows from the exponent $\exp(-6976/T)$ of their
Eq.~(14), giving $E_a = 6976R = \SI{58.0}{\kilo\joule\per\mole}$ --- not
the \SI{24}{\kilo\joule\per\mole} used in an earlier version of this
formulation, which understates the temperature sensitivity by a factor
of \num{2.42} and the \SIrange{25}{40}{\celsius} acceleration by
\num{1.93}. The cycle throughput exponent $\beta = 0.5$ is their
Eq.~(18). And the depth coefficient is fixed by the ratio of the two
depth terms in their Eq.~(20),
$c_D = \num{4.081e-3} / \num{7.600e-4} = \num{5.3697}$, against the
value of unity previously assumed; the ratio is invariant to whether
their charge throughput counts one or both directions.

Four things are \emph{not} theirs and must not be attributed. Their
calendar ageing is parameterised on cell \emph{voltage}, linearly,
$a_{\mathrm{cap}} \propto 7.543V - 23.75$; converting to the
state-of-charge form of \eqref{eq:gstress} requires the OCV curve their
Fig.~1 plots but does not tabulate, so the hinge shape here remains a
surrogate. Their depth term is \emph{additive} where
\eqref{eq:cyc} is multiplicative; the two agree at the fitted baseline
and diverge away from it. Every cycle test ran at 1C and
\SI{35}{\celsius}, so the study supports no C-rate and no
cycle-temperature dependence at all --- $c_C$ is therefore set to zero
rather than invented, and treated as a sensitivity parameter alongside
$\lambda_d$. Finally their Eq.~(20) carries a mean-voltage term,
$\num{7.348e-3}(\bar{V} - 3.667)^2$, with no counterpart here; it is
not small, since cycling about \SI{4.0}{\volt} rather than
\SI{3.667}{\volt} doubles their cycle coefficient. Omitting it
understates cycling at high mean state of charge, which is where V2G
operates.

\paragraph{Purewal et al.\ (2014) is not the NMC source.}
An earlier version of this work cited ``Schmalstieg et al.\ (2014)''
against the title of Wang/Purewal et al.~\cite{wang2014,purewal2014},
conflating two distinct 2014 \emph{J.\ Power Sources} papers. Purewal
et al.\ is a physics-based chemical--mechanical model at fixed
\SI{50}{\percent} depth of discharge, with no state-of-charge stress
function, a $\sqrt{t}$ calendar term, and cycle ageing as Paris'-law
fatigue crack growth ($m = 2.2$). None of it is used here. It is
retained in the bibliography because its Part~1 companion reports the
empirical activation energy \SI{24.5}{\kilo\joule\per\mole} that the
superseded \SI{24}{\kilo\joule\per\mole} most plausibly came from.

\paragraph{Pack costs.}
$\pi_{\mathrm{LFP}}$ and $\pi_{\mathrm{NMC}}$ derive from a reported
chemistry split of \SI{81}{\USD\per\kilo\watt\hour} and
\SI{128}{\USD\per\kilo\watt\hour} at approximately
\SI{96.81}{\rupee\per\USD}, giving \SI{7841.6}{\rupee} and
\SI{12391.7}{\rupee}, rounded to the table values. Three cautions. These
are \emph{global volume-weighted pack prices across applications}, not
Indian EV replacement-pack retail quotes; the same source separately
reports BEV packs averaging \SI{99}{\USD\per\kilo\watt\hour}, equivalent
to \SI{9584}{\rupee\per\kilo\watt\hour}, so the NMC figure sits about
\SI{29}{\percent} above the BEV-specific average. The exchange rate is
undated. And this is the single exogenous monetary assumption in the
model, so the recommended $\pm\SI{25}{\percent}$ band should accompany
any quoted rupee figure.

\paragraph{Temperature.}
Equation~\eqref{eq:arrhenius} propagates strongly. Pack temperature is
held at a scenario constant here; fade precision beyond what that
assumption supports should not be quoted.

\paragraph{Multi-day coupling.}
$N$ is a scenario input and is not advanced between sessions in the
charge-only results. This is benign where $\varphi$ is small and
$\sigma$ slowly varying, and degrades under V2G, where per-session
throughput is several times larger.

\paragraph{Asymmetric regularisation of the two power laws.}
The cycle slope \eqref{eq:cycslope} is floored at $N^{\min}$, bounding
$\sigma \le \beta (N^{\min})^{\beta-1} = \num{0.5}$ for a pack with no
accumulated throughput. The age factor \eqref{eq:agefactor} carries no
counterpart: $\psi$ grows without bound as $t \to 0$
($\psi = \num{1.78}$ at \SI{0.1}{\yr}, \num{3.16} at \SI{0.01}{\yr} for
NMC) and only a non-positive age is rejected. Both expressions are
singular derivatives of the same power-law form, so the absence of a
floor on one is an inconsistency rather than a modelling choice. Results
for near-new packs should be treated with care until an explicit minimum
anchor age is introduced.

\paragraph{Depth measure.}
Equation~\eqref{eq:delta} uses the full-session excursion as the depth
proxy. Validation against rainflow counting is outstanding and should
compare \emph{offer rankings} rather than absolute fade.

\begin{thebibliography}{9}

\bibitem{naumann2018}
M.~Naumann, M.~Schimpe, P.~Keil, H.~C. Hesse, A.~Jossen,
``Analysis and modeling of calendar aging of a commercial
LiFePO$_4$/graphite cell,''
\emph{Journal of Energy Storage}, vol.~17, pp.~153--169, 2018.
doi:10.1016/j.est.2018.01.019

\bibitem{wang2014}
J.~Wang, J.~Purewal, P.~Liu, J.~Hicks-Garner, S.~Soukazian, E.~Sherman,
A.~Sorenson, L.~Vu, H.~Tataria, M.~W. Verbrugge,
``Degradation of lithium ion batteries employing graphite negatives and
nickel-cobalt-manganese oxide $+$ spinel manganese oxide positives:
Part~1, aging mechanisms and life estimation,''
\emph{Journal of Power Sources}, vol.~269, pp.~937--948, 2014.
doi:10.1016/j.jpowsour.2014.07.030

\bibitem{purewal2014}
J.~Purewal, J.~Wang, J.~Graetz, S.~Soukiazian, H.~Tataria,
M.~W. Verbrugge,
``Degradation of lithium ion batteries employing graphite negatives and
nickel-cobalt-manganese $+$ spinel manganese oxide positives:
Part~2, chemical-mechanical degradation model,''
\emph{Journal of Power Sources}, 2014.

\bibitem{schmalstieg2014}
J.~Schmalstieg, S.~K\"abitz, M.~Ecker, D.~U. Sauer,
``A holistic aging model for Li(NiMnCo)O$_2$ based 18650 lithium-ion
batteries,''
\emph{Journal of Power Sources}, vol.~257, pp.~325--334, 2014.
doi:10.1016/j.jpowsour.2014.02.012

\end{thebibliography}

\end{document}
